% from ega4-14-fr.tex
\oldpage[IV-3]{202}
\subsection{Morphismes ouverts et formule des dimensions}
\label{subsection:IV.14.2-fr}

\begin{theorem}[14.2.1]
\EGAReaderTarget{ega4.c04.thm.14-2-1}{IV.14.2.1-fr}
Soient $X$, $Y$ deux préschémas localement noethériens, $f:X\to Y$ un morphisme,
$x$ un point de $X$, $y=f(x)$. On suppose que $f$ est ouvert aux points génériques
des composantes irréductibles de $f^{-1}(y)$ contenant $x$. Alors on a la relation
\begin{equation}
\tag{14.2.1.1}
\label{IV.14.2.1.1-fr}
  \dim(\mathcal O_x)=\dim(\mathcal O_y)+
  \dim(\mathcal O_x\otimes_{\mathcal O_y}k(y)).
\end{equation}
\end{theorem}

\begin{proof}
Soit $y'$ une générisation quelconque de $y$ distincte de $y$, et considérons le
sous-préschéma fermé réduit $Y'$ de $Y$ ayant pour espace sous-jacent
$\overline{\{y'\}}$~; alors aucune composante irréductible $X'$ de $f^{-1}(Y')$
contenant $x$ ne peut être contenue dans $f^{-1}(y)$~: en effet, si $x'$ est le
point générique de $X'$, on aurait $f(x')=y\ne y'$, et comme $x'$ est sa seule
générisation dans $f^{-1}(Y')$, cela contredirait l'hypothèse que $f$ est ouvert
au point $x'$, en vertu du fait que $a)$ implique $c)$ dans (1.10.3). On peut donc
appliquer (6.1.2) à l'homomorphisme local d'anneaux noethériens
$\mathcal O_y\to\mathcal O_x$, d'où la conclusion.
\end{proof}

\begin{corollary}[14.2.2]
\label{IV.14.2.2-fr}
Soient $Y$ un préschéma localement noethérien, $f:X\to Y$ un morphisme localement
de type fini, $x$ un point de $X$, $y=f(x)$. Supposons $f$ ouvert aux points
génériques des composantes irréductibles de $f^{-1}(y)$ contenant $x$. Alors $f$
est équidimensionnel au point $x$ dans chacun des deux cas suivants~:
\begin{enumerate}[label=(\roman*)]
  \item $X$ est irréductible et $f$ est dominant.
  \item Les anneaux $\mathcal O_x$ et $\mathcal O_y$ sont équidimensionnels.
\end{enumerate}
\end{corollary}

\begin{proof}
Pour (i), cela résulte de (14.2.1) et de (13.2.3). Pour (ii), cela résulte de
(14.2.1) et de (13.3.6).
\end{proof}

\begin{proposition}[14.2.3]
\label{IV.14.2.3-fr}
Soient $Y$ un préschéma irréductible localement noethérien, $\eta$ son point
générique, $f:X\to Y$ un morphisme localement de type fini, $y$ un point de
$f(X)$, $Z$ une composante irréductible de $f^{-1}(y)$ telle que $f$ soit ouvert
au point générique de $Z$. Alors $Z$ est contenue dans une composante irréductible
$X'$ de $X$ dominant $Y$ et telle que
\[
  \dim(Z)=\dim(X'\cap f^{-1}(y))=\dim(X'\cap f^{-1}(\eta)).
\]
\end{proposition}

\begin{proof}
Cela résulte en effet de (14.2.1) appliqué au point générique de $Z$, et de
(13.2.7).
\end{proof}

\begin{corollary}[14.2.4]
\label{IV.14.2.4-fr}
Avec les notations de (14.2.3), supposons que $f$ soit ouvert en tous les points
de $f^{-1}(y)$. Pour tout $z\in Y$, soient $E(z)$ l'ensemble des dimensions des
composantes irréductibles de $f^{-1}(z)$ et
$d(z)=\sup E(z)=\dim(f^{-1}(z))$. On a alors $E(y)\subset E(\eta)$, d'où
$d(y)\leq d(\eta)$.
\end{corollary}

\begin{proof}
Cela résulte aussitôt de (14.2.3).
\end{proof}

\begin{corollary}[14.2.5]
\label{IV.14.2.5-fr}
Soient $Y$ un préschéma localement noethérien, $f:X\to Y$ un morphisme propre,
$y\in Y$ un point tel que $f$ soit ouvert en tous les points de $f^{-1}(y)$. Alors
la fonction $z\mapsto\dim(f^{-1}(z))$ est constante dans un voisinage de $y$.
\end{corollary}

\begin{proof}
Soient $Y_i$ ($1\leq i\leq n$) les sous-préschémas fermés réduits de $Y$ ayant
pour espaces sous-jacents les composantes irréductibles de $Y$ contenant $y$~;
si $f_i:f^{-1}(Y_i)\to Y_i$ est la restriction de $f$, on sait que chacun des
$f_i$ est propre (\textbf{II}, 5.4.5) et ouvert en tous les points de $f^{-1}(y)$
(14.1.2). Comme la réunion des $Y_i$ est un voisinage de $y$ dans $Y$,
